\documentclass[12pt]{article}
\usepackage{amsmath, amssymb, amsthm}
\usepackage{geometry}
\geometry{a4paper, margin=1in}
\usepackage{enumitem}
\usepackage{parskip}

\title{Homework 16.3: Fundamental Theorem of Line Integrals \\ \Large Problem Sets and Solutions}
\author{}
\date{}

\begin{document}

\maketitle

\section*{Part 1: 2D Conservative Field Verification}

\textbf{Problem 1:} Determine if $\mathbf{F}(x,y) = \langle 2xy^3, 3x^2y^2 \rangle$ is a conservative vector field.
\textbf{Solution:} Let $P = 2xy^3$ and $Q = 3x^2y^2$. 
Compute the cross-partials: $P_y = \frac{\partial}{\partial y}(2xy^3) = 6xy^2$ and $Q_x = \frac{\partial}{\partial x}(3x^2y^2) = 6xy^2$. 
Since $P_y = Q_x$, $\mathbf{F}$ is conservative.

\textbf{Problem 2:} Determine if $\mathbf{F}(x,y) = \langle \sin(y), \cos(x) \rangle$ is a conservative vector field.
\textbf{Solution:} Let $P = \sin(y)$ and $Q = \cos(x)$. 
Compute $P_y = \cos(y)$ and $Q_x = -\sin(x)$. 
Since $\cos(y) \neq -\sin(x)$, $\mathbf{F}$ is not conservative.

\textbf{Problem 3:} Determine if $\mathbf{F}(x,y) = \langle e^x \sin(y), e^x \cos(y) \rangle$ is a conservative vector field.
\textbf{Solution:} Let $P = e^x \sin(y)$ and $Q = e^x \cos(y)$. 
Compute $P_y = e^x \cos(y)$ and $Q_x = e^x \cos(y)$. 
Since $P_y = Q_x$, $\mathbf{F}$ is conservative.

\textbf{Problem 4:} Determine if $\mathbf{F}(x,y) = \langle 4x^3 y + 2x, x^4 - y^2 \rangle$ is a conservative vector field.
\textbf{Solution:} Let $P = 4x^3 y + 2x$ and $Q = x^4 - y^2$. 
Compute $P_y = 4x^3$ and $Q_x = 4x^3$. 
Since $P_y = Q_x$, $\mathbf{F}$ is conservative.

\textbf{Problem 5:} Determine if $\mathbf{F}(x,y) = \left\langle \frac{x}{x^2+y^2}, \frac{y}{x^2+y^2} \right\rangle$ is a conservative vector field.
\textbf{Solution:} Let $P = \frac{x}{x^2+y^2}$ and $Q = \frac{y}{x^2+y^2}$. 
Using the quotient rule: 
$P_y = \frac{0 - x(2y)}{(x^2+y^2)^2} = \frac{-2xy}{(x^2+y^2)^2}$ and 
$Q_x = \frac{0 - y(2x)}{(x^2+y^2)^2} = \frac{-2xy}{(x^2+y^2)^2}$. 
Since $P_y = Q_x$, $\mathbf{F}$ is conservative.

\section*{Part 2: 3D Conservative Field Verification}

\textbf{Problem 6:} Determine if $\mathbf{F}(x,y,z) = \langle yz^2, xz^2, 2xyz \rangle$ is a conservative vector field.
\textbf{Solution:} We check if the curl $\nabla \times \mathbf{F} = \mathbf{0}$ by testing the 3D components $P, Q, R$.
1) $R_y = 2xz$ and $Q_z = 2xz \implies R_y = Q_z$.
2) $P_z = 2yz$ and $R_x = 2yz \implies P_z = R_x$.
3) $Q_x = z^2$ and $P_y = z^2 \implies Q_x = P_y$.
Since all cross-partials match, the curl is zero, and $\mathbf{F}$ is conservative.

\textbf{Problem 7:} Determine if $\mathbf{F}(x,y,z) = \langle y, z, x \rangle$ is a conservative vector field.
\textbf{Solution:} Check $R_y$ and $Q_z$. 
$R = x \implies R_y = 0$. $Q = z \implies Q_z = 1$. 
Since $0 \neq 1$, the field is not conservative.

\textbf{Problem 8:} Determine if $\mathbf{F}(x,y,z) = \langle e^z, 1, x e^z \rangle$ is a conservative vector field.
\textbf{Solution:} 
1) $R_y = \frac{\partial}{\partial y}(x e^z) = 0$ and $Q_z = \frac{\partial}{\partial z}(1) = 0$.
2) $P_z = \frac{\partial}{\partial z}(e^z) = e^z$ and $R_x = \frac{\partial}{\partial x}(x e^z) = e^z$.
3) $Q_x = \frac{\partial}{\partial x}(1) = 0$ and $P_y = \frac{\partial}{\partial y}(e^z) = 0$.
All cross-partials match, so $\mathbf{F}$ is conservative.

\section*{Part 3: Finding 2D Scalar Potential Functions}

\textbf{Problem 9:} Find a potential function $f(x,y)$ for $\mathbf{F}(x,y) = \langle 6xy^2 - 5, 6x^2y + 2y \rangle$.
\textbf{Solution:} 
1. Integrate $P$ with respect to $x$: $f(x,y) = \int (6xy^2 - 5) dx = 3x^2y^2 - 5x + g(y)$.
2. Differentiate $f$ w.r.t $y$ and equate to $Q$: $f_y = 6x^2y + g'(y) = 6x^2y + 2y$.
3. Solve for $g(y)$: $g'(y) = 2y \implies g(y) = y^2$. 
\textbf{Final Answer:} $f(x,y) = 3x^2y^2 - 5x + y^2$.

\textbf{Problem 10:} Find a potential function $f(x,y)$ for $\mathbf{F}(x,y) = \langle y e^{xy} + 2x, x e^{xy} - 3y^2 \rangle$.
\textbf{Solution:} 
1. $f(x,y) = \int (y e^{xy} + 2x) dx = e^{xy} + x^2 + g(y)$.
2. $f_y = x e^{xy} + g'(y) = x e^{xy} - 3y^2$.
3. $g'(y) = -3y^2 \implies g(y) = -y^3$.
\textbf{Final Answer:} $f(x,y) = e^{xy} + x^2 - y^3$.

\textbf{Problem 11:} Find a potential function $f(x,y)$ for the rational field $\mathbf{F}(x,y) = \langle \frac{1}{x} + y, x + \frac{1}{y} \rangle$ for $x,y > 0$.
\textbf{Solution:} 
1. $f = \int (\frac{1}{x} + y) dx = \ln(x) + xy + g(y)$.
2. $f_y = x + g'(y) = x + \frac{1}{y}$.
3. $g'(y) = \frac{1}{y} \implies g(y) = \ln(y)$.
\textbf{Final Answer:} $f(x,y) = \ln(x) + \ln(y) + xy$.

\textbf{Problem 12:} Find a potential function $f(x,y)$ for $\mathbf{F}(x,y) = \langle 2x \cos(y), -x^2 \sin(y) + 4 \rangle$.
\textbf{Solution:} 
1. $f = \int 2x \cos(y) dx = x^2 \cos(y) + g(y)$.
2. $f_y = -x^2 \sin(y) + g'(y) = -x^2 \sin(y) + 4$.
3. $g'(y) = 4 \implies g(y) = 4y$.
\textbf{Final Answer:} $f(x,y) = x^2 \cos(y) + 4y$.

\textbf{Problem 13:} Find a potential function $f(x,y)$ for $\mathbf{F}(x,y) = \langle y^2 \sin(xy^2), 2xy \sin(xy^2) \rangle$.
\textbf{Solution:} 
1. Use integration by substitution. Let $u = xy^2$, $du = y^2 dx$. 
$f = \int y^2 \sin(xy^2) dx = -\cos(xy^2) + g(y)$.
2. $f_y = -\left(-\sin(xy^2) \cdot 2xy\right) + g'(y) = 2xy \sin(xy^2) + g'(y)$.
Equating to $Q$: $2xy \sin(xy^2) + g'(y) = 2xy \sin(xy^2) \implies g'(y) = 0$.
\textbf{Final Answer:} $f(x,y) = -\cos(xy^2)$.

\section*{Part 4: Finding 3D Scalar Potential Functions}

\textbf{Problem 14:} Find a potential function $f(x,y,z)$ for $\mathbf{F}(x,y,z) = \langle 2xy, x^2 + 2yz, y^2 \rangle$.
\textbf{Solution:} 
1. Integrate $P$ w.r.t $x$: $f = \int 2xy dx = x^2y + g(y,z)$.
2. Differentiate w.r.t $y$ and equate to $Q$: $f_y = x^2 + g_y(y,z) = x^2 + 2yz \implies g_y(y,z) = 2yz$.
3. Integrate $g_y$ w.r.t $y$: $g(y,z) = y^2z + h(z) \implies f = x^2y + y^2z + h(z)$.
4. Differentiate w.r.t $z$ and equate to $R$: $f_z = y^2 + h'(z) = y^2 \implies h'(z) = 0 \implies h(z) = K$.
\textbf{Final Answer:} $f(x,y,z) = x^2y + y^2z$.

\textbf{Problem 15:} Find a potential function $f(x,y,z)$ for $\mathbf{F}(x,y,z) = \langle \ln(y), \frac{x}{y} + z, y \rangle$ for $y>0$.
\textbf{Solution:} 
1. $f = \int \ln(y) dx = x\ln(y) + g(y,z)$.
2. $f_y = \frac{x}{y} + g_y = \frac{x}{y} + z \implies g_y = z \implies g(y,z) = yz + h(z)$. 
$f = x\ln(y) + yz + h(z)$.
3. $f_z = y + h'(z) = y \implies h'(z) = 0$.
\textbf{Final Answer:} $f(x,y,z) = x\ln(y) + yz$.

\textbf{Problem 16:} Find a potential function $f(x,y,z)$ for $\mathbf{F}(x,y,z) = \langle z e^{xz}, 2y, x e^{xz} + 3z^2 \rangle$.
\textbf{Solution:} 
1. $f = \int z e^{xz} dx = e^{xz} + g(y,z)$.
2. $f_y = g_y = 2y \implies g(y,z) = y^2 + h(z)$. $f = e^{xz} + y^2 + h(z)$.
3. $f_z = x e^{xz} + h'(z) = x e^{xz} + 3z^2 \implies h'(z) = 3z^2 \implies h(z) = z^3$.
\textbf{Final Answer:} $f(x,y,z) = e^{xz} + y^2 + z^3$.

\section*{Part 5: Evaluating Line Integrals using FTLI in 2D}

\textbf{Problem 17:} Evaluate $\int_C \langle 3x^2y^2, 2x^3y \rangle \cdot d\mathbf{r}$ along the curve $C: y = x^{3/2}$ from $(0,0)$ to $(1,1)$.
\textbf{Solution:} The field is conservative with potential $f = x^3y^2$. Using the FTLI, we evaluate at the endpoints: 
$\int_C \mathbf{F} \cdot d\mathbf{r} = f(1,1) - f(0,0) = (1^3 \cdot 1^2) - (0) = 1$.

\textbf{Problem 18:} Evaluate $\int_C \langle e^y, x e^y + 2y \rangle \cdot d\mathbf{r}$ along the arc of the unit circle in the first quadrant from $(1,0)$ to $(0,1)$.
\textbf{Solution:} Since evaluating the path directly is messy, we find a potential: $f = x e^y + y^2$.
Using FTLI: $f(0,1) - f(1,0) = (0 \cdot e^1 + 1^2) - (1 \cdot e^0 + 0^2) = 1 - 1 = 0$.

\textbf{Problem 19:} Evaluate $\int_C \langle y \cos(xy), x \cos(xy) \rangle \cdot d\mathbf{r}$ from $(\pi, 1)$ to $(1, \pi/2)$.
\textbf{Solution:} The potential function is $f(x,y) = \sin(xy)$. 
$f(1, \pi/2) - f(\pi, 1) = \sin(\pi/2) - \sin(\pi) = 1 - 0 = 1$.

\textbf{Problem 20:} Evaluate $\int_C \langle \frac{y^2}{x}, 2y \ln(x) \rangle \cdot d\mathbf{r}$ from $(1, 1)$ to $(e, 2)$.
\textbf{Solution:} The potential function is $f(x,y) = y^2 \ln(x)$. 
$f(e, 2) - f(1, 1) = (2^2 \ln(e)) - (1^2 \ln(1)) = 4(1) - 1(0) = 4$.

\textbf{Problem 21:} Evaluate $\int_C \langle 2x - y, -x + 4y^3 \rangle \cdot d\mathbf{r}$ from $(2, 0)$ to $(0, 2)$.
\textbf{Solution:} $f = \int(2x-y)dx = x^2 - xy + g(y)$. 
$f_y = -x + g'(y) = -x + 4y^3 \implies g(y) = y^4$. Thus, $f(x,y) = x^2 - xy + y^4$.
$f(0, 2) - f(2, 0) = (0 - 0 + 16) - (4 - 0 + 0) = 16 - 4 = 12$.

\section*{Part 6: Evaluating Line Integrals using FTLI in 3D}

\textbf{Problem 22:} Evaluate $\int_C \langle yz, xz, xy \rangle \cdot d\mathbf{r}$ from $(1, 1, 1)$ to $(2, 3, 4)$.
\textbf{Solution:} The potential function is $f(x,y,z) = xyz$. 
$\int_C \mathbf{F} \cdot d\mathbf{r} = f(2,3,4) - f(1,1,1) = (2)(3)(4) - (1)(1)(1) = 24 - 1 = 23$.

\textbf{Problem 23:} Evaluate $\int_C \langle 2x e^z, 2y, x^2 e^z \rangle \cdot d\mathbf{r}$ from $(1, 2, 0)$ to $(3, 0, \ln(2))$.
\textbf{Solution:} 
Potential $f = x^2 e^z + y^2$. 
$f(3, 0, \ln 2) - f(1, 2, 0) = \left(9 e^{\ln 2} + 0\right) - \left(1 e^0 + 4\right) = 9(2) - (1 + 4) = 18 - 5 = 13$.

\textbf{Problem 24:} Evaluate $\int_C \langle \sin(y), x \cos(y) + z, y \rangle \cdot d\mathbf{r}$ from $(0, 0, 0)$ to $(2, \pi/2, 3)$.
\textbf{Solution:} 
1. $f = \int \sin(y) dx = x \sin(y) + g(y,z)$.
2. $f_y = x \cos(y) + g_y = x \cos(y) + z \implies g_y = z \implies g = yz + h(z)$.
3. $f_z = y + h'(z) = y \implies h'(z)=0 \implies f = x \sin(y) + yz$.
$f(2, \pi/2, 3) - f(0,0,0) = \left(2\sin(\pi/2) + (\pi/2)(3)\right) - 0 = 2 + \frac{3\pi}{2}$.

\section*{Part 7: The Closed Loop Property}

\textbf{Problem 25:} Evaluate $\int_C \langle 4x^3, 3y^2 \rangle \cdot d\mathbf{r}$ where $C$ is the closed ellipse $\frac{x^2}{4} + \frac{y^2}{9} = 1$.
\textbf{Solution:} The field is conservative because $P_y = 0$ and $Q_x = 0$. Because $C$ is a closed loop, the start and end points are identical. The FTLI guarantees $\int_C \mathbf{F} \cdot d\mathbf{r} = 0$.

\textbf{Problem 26:} Evaluate $\int_C \langle z e^{xz}, 2y, x e^{xz} \rangle \cdot d\mathbf{r}$ where $C$ is the boundary of the rectangle $0 \le x \le 1, 0 \le y \le 2$ at the plane $z=1$.
\textbf{Solution:} Checking curl reveals $\nabla \times \mathbf{F} = \mathbf{0}$, meaning the field is conservative. Because $C$ is a boundary (a closed loop), the line integral is exactly $0$.

\section*{Part 8: Real-World Applications (Finance Focus)}

\textbf{Problem 27:} An arbitrage pricing model features a local cash sensitivity vector field $\mathbf{F}(x,y) = \langle 2xy^3, 3x^2y^2 \rangle$. Verify the market no-arbitrage condition (verify the field is conservative) and find the absolute pricing potential function $V(x,y)$ such that $\nabla V = \mathbf{F}$.
\textbf{Solution:} 
1. $P_y = 6xy^2$ and $Q_x = 6xy^2$. Since they match, the market is arbitrage-free.
2. $V(x,y) = \int 2xy^3 dx = x^2y^3 + g(y)$. 
3. $V_y = 3x^2y^2 + g'(y) = 3x^2y^2 \implies g'(y)=0$. 
\textbf{Final Answer:} $V(x,y) = x^2y^3$.

\textbf{Problem 28:} An automated trading algorithm trades based on the vector field of required cash flows $\mathbf{F}(S_1, S_2) = \langle e^{S_2}, S_1 e^{S_2} \rangle$. It traces a path $C$ from Monday's opening prices $(1, 0)$ to Friday's closing prices $(5, \ln 4)$. Calculate the total cost of trading over the week.
\textbf{Solution:} 
The field is conservative with potential $V(S_1, S_2) = S_1 e^{S_2}$. 
The total cost is independent of the path. 
$Cost = V(5, \ln 4) - V(1, 0) = 5e^{\ln 4} - 1e^0 = 5(4) - 1 = 19$.

\section*{Part 9: Advanced Diagnostic Testing (Find the Flaw)}

\textbf{Problem 29:} 
\textbf{Problem:} Evaluate $\int_C \langle x^2y, xy^2 \rangle \cdot d\mathbf{r}$ along a curve from $(0,0)$ to $(1,1)$. 
\textbf{Flawed Solution:} We use the FTLI. $f_x = x^2y \implies f = \frac{1}{3}x^3y$. So $\int_C \mathbf{F} \cdot d\mathbf{r} = \frac{1}{3}(1)^3(1) - 0 = 1/3$.
\textbf{Spot the Flaw:} Explain why the solution is conceptually wrong and provide the correct reasoning.
\textbf{Solution:} The student failed to test if the field is conservative first. $P_y = x^2$ and $Q_x = y^2$. Since $x^2 \neq y^2$, the field is NOT conservative. The FTLI cannot be used; the integral depends entirely on the specific path taken (which requires parameterization).

\textbf{Problem 30:} 
\textbf{Problem:} Find a potential function for $\mathbf{F} = \langle 2x \cos(y), -x^2 \sin(y) + 2y \rangle$. 
\textbf{Flawed Solution:} 1. $f_x = 2x \cos(y) \implies f = x^2 \cos(y) + C$. 2. We check $f_y = -x^2 \sin(y)$, which does not match $Q = -x^2 \sin(y) + 2y$. 3. Therefore, the field is not conservative and no potential exists.
\textbf{Spot the Flaw:} Identify the conceptual error and find the true potential function.
\textbf{Solution:} The student integrated with respect to $x$ and added a numerical constant $+C$ instead of an arbitrary function $+g(y)$. This is a critical conceptual misunderstanding of partial integration. If they had used $+g(y)$, they would have found $g'(y) = 2y \implies g(y) = y^2$, correctly yielding the potential $f(x,y) = x^2 \cos(y) + y^2$.

\newpage

\section*{Concept Check List}
Below is a mapping of the core concepts tested across all 30 problems to easily target specific skills.

\begin{itemize}[label=$\square$]
    \item \textbf{[C1] 2D Conservative Field Check}: Understanding $P_y = Q_x$. \newline \textit{Relevant Problems: 1, 2, 3, 4, 5}
    \item \textbf{[C2] 3D Conservative Field Check}: Applying Curl $\nabla \times \mathbf{F} = \mathbf{0}$. \newline \textit{Relevant Problems: 6, 7, 8}
    \item \textbf{[C3] Finding 2D Scalar Potentials}: Polynomials, transcendentals, rationals, and treating variables as constants (+g(y)). \newline \textit{Relevant Problems: 9, 10, 11, 12, 13}
    \item \textbf{[C4] Finding 3D Scalar Potentials}: Triple integration-differentiation loop. \newline \textit{Relevant Problems: 14, 15, 16}
    \item \textbf{[C5] Evaluating Line Integrals using FTLI (2D)}: Skipping path parameterization for $f(B) - f(A)$. \newline \textit{Relevant Problems: 17, 18, 19, 20, 21}
    \item \textbf{[C6] Evaluating Line Integrals using FTLI (3D)}: Extending the FTLI to 3D endpoints. \newline \textit{Relevant Problems: 22, 23, 24}
    \item \textbf{[C7] Closed Loop Property}: Identifying loops and applying $\int_C = 0$. \newline \textit{Relevant Problems: 25, 26}
    \item \textbf{[C8] Real-World Application}: Translating finance terminology (cash sensitivities, arbitrage) to gradient math. \newline \textit{Relevant Problems: 27, 28}
    \item \textbf{[C9] Conceptual Debugging (Find the Flaw)}: Spotting blind FTLI applications and forgotten $+g(y)$ constants. \newline \textit{Relevant Problems: 29, 30}
\end{itemize}

\end{document}